% !PRINTED 347
\weberpage{347}

%<<CONT
and determine the sign so that $\pm q \equiv 1 \ (\mathrm{mod}\ 4)$.
In this case
\[
\theta = \frac{-1 + \sqrt{\pm q}}{2},
\]
and we can represent this expression by $q$th roots of unity.
Let $r$ denote an imaginary $q$th root of unity, and let
$a$ run through the quadratic residues, $b$ through the non-residues of $q$.
We set, as in Bd.\ I, \S\,179
\weq{4}{A = \overset{a}{\Sigma}\, r^a, \quad B = \overset{b}{\Sigma}\, r^b}
and have, by the formulas (3) given there, which are derived without the use
of the law of reciprocity:
\weq{5}{\theta = A, \quad \theta' = B.}
The sign of $\sqrt{+q}$ depends on the choice of $r$,
but does not come into consideration here. Now, by the polynomial
 theorem:
\[
A^p \equiv \Sigma r^{pa}, \quad B^p \equiv \Sigma r^{pb},
\]
hence
\[
A^p \equiv A, \quad B^p \equiv B, \quad \text{if} \ \left(\frac{p}{q}\right) = +1,
\]
and
\[
A^p \equiv B, \quad B^p \equiv A, \quad \text{if} \ \left(\frac{p}{q}\right) = -1.
\]

Since $A$ is not congruent to $B$ when $p$ is different from $q$,
it follows from this and from 2. of the preceding section:

3) \gesperrt{The congruence}
\[
x^2 \equiv \pm q \ (\mathrm{mod}\ p)
\]
\gesperrt{has solutions then and only then, when $p$ is a quadratic
residue of $q$.}

And this is the law of reciprocity with its supplementary theorems.

\S\,95. \gesperrt{Equivalent Forms and Ideals in the Quadratic Field.}

If the two ideals [\S\,91, (16)]
\weq{1}{\mathfrak{a} = (a,\, a\,\omega), \quad \mathfrak{a}' = (a',\, a'\,\omega')}
are equivalent, then there is an (integral or non-integral) number $\eta$
in $\Omega$ for which $\eta\, \mathfrak{a}' = \mathfrak{a}$, and this number $\eta$ is determined by the ideals
$\mathfrak{a}$, $\mathfrak{a}'$ themselves up to a factor which is a \gesperrt{numerical unit}
%CONT>>
